\section*{第 \ref{ch:b2:sound-waves} 章——流体中的声波}

\begin{solution}{exo:b2:sound-waves:1}
$c = \sqrt{\gamma RT/M}$：\qty{331}{m/s}、\qty{343}{m/s}、\qty{299}{m/s}；氦气中
\qty{1010}{m/s}；CO$_2$ 中 \qty{268}{m/s}。声带以相同
频率振动，但声道的共振（决定音色）与 $c$ 成正比，
故升高为原来的约三倍：
声音听起来音调很高。
\end{solution}

\begin{solution}{exo:b2:sound-waves:2}
\qty{60}{dB}: $I = \qty{1e-6}{W/m^2}$, $p_m = \sqrt{2ZI} = \qty{29}{mPa}$, $v_m = p_m/Z =
\qty{7e-5}{m/s}$, $\xi_m = v_m/\omega = \qty{22}{nm}$. \qty{120}{dB}: \qty{1}{W/m^2},
\qty{29}{Pa}，\qty{7}{cm/s}，\qty{22}{\micro m}。两个声源：\qty{63}{dB}；一百个声源：\qty{80}{dB}。
\end{solution}

\begin{solution}{exo:b2:sound-waves:3}
$I = 1/4\pi r^2$: \qty{0.080}{W/m^2} (\qty{109}{dB}), \qty{89}{dB}, \qty{69}{dB}; $I = 10^{-6}$
$r = \sqrt{1/4\pi10^{-6}} = \qty{280}{m}$ 处；\qty{100}{dB} 在 \qty{30}{m} 处需要 $4\pi \times
900 \times 10^{-2} = \qty{110}{W}$ 的声功率。
\end{solution}

\begin{solution}{exo:b2:sound-waves:4}
$4.5 \times 340 = \qty{1.5}{km}$; $340 \times 1.5 = \qty{510}{m}$; $1500 \times 1.2 = \qty{1.8}{km}$;
$1540 \times 65 \times 10^{-6} = \qty{10}{cm}$.
\end{solution}

\begin{solution}{exo:b2:sound-waves:5}
(a) $500 \times 340/300 = \qty{567}{Hz}$；$500 \times 340/380 = \qty{447}{Hz}$。(b) $500 \times
380/340 = \qty{559}{Hz}$；两者同时运动：$500 \times 380/300 = \qty{633}{Hz}$，
与仅声源运动（速度 \qty{80}{m/s}）时的 \qty{654}{Hz}，或仅观察者运动时的 \qty{618}{Hz} 相比：起决定作用的是空气，
而不是相对速度。(c) 墙壁接收到 \qty{567}{Hz}
的声波，并静止地再发出；驱动者以 \qty{40}{m/s} 的速度靠近，频率为 $567 \times
380/340 = \qty{633}{Hz}$.
\end{solution}

\begin{solution}{exo:b2:sound-waves:6}
$v_m = p_m/Z$: \qty{2.4}{mm/s}, \qty{0.67}{\micro m/s}, \qty{25}{nm/s}; $\xi_m = v_m/\omega$:
\qty{390}{nm}, \qty{0.11}{nm}, \qty{4}{pm}; $I = p_m^2/2Z$: \qty{1.2e-3}{W/m^2} (\qty{91}{dB}),
\qty{3.3e-7}{W/m^2}（\qty{55}{dB}）、\qty{1.3e-8}{W/m^2}（\qty{41}{dB}）。刚性介质
在给定压强下几乎不动，并传递很小的功率。
\qty{1}{Pa} 相当于 \qty{120}{dB}（参考值为 \qty{1}{\micro Pa}）。
\end{solution}

\begin{solution}{exo:b2:sound-waves:7}
(a) $\sqrt{1.013 \times 10^5/1.2} = \qty{290}{m/s}$；$\times\sqrt{1.4}$：\qty{343}{m/s}。(b)
$\gamma = c^2M/RT$：空气为 $1.40$，氩气为 $1.71$（在取整误差内为 $5/3$）：
三个平动自由度（$\gamma = 1 + 2/3$），
而双原子气体为五个（$1 + 2/5$）。(c) $\sqrt{D/f} = \sqrt{2 \times 10^{-8}} = \qty{0.14}{mm}
\ll \lambda = \qty{34}{cm}$：波峰与波谷之间没有热量流动。
\end{solution}

\begin{solution}{exo:b2:sound-waves:8}
(a) 两端开口：$p = p_m\sin kx\cos\omega t$，$kL = n\pi$；一端闭口、一端开口：
$p = p_m\cos kx\cos\omega t$（速度节、压强腹位于闭口端），
$kL = (2n - 1)\pi/2$。(b) $L = c/2f = \qty{39}{cm}$；$c/4f = \qty{19.5}{cm}$。
(c) 全部谐音都会出现；仅奇次谐音（即单簧管特有的空洞声音）。
(d) 每个开口端的修正量为 $0.6 \times 2 = \qty{1.2}{cm}$：$L_{\text{有效}} = \qty{41.4}{cm}$，$f =
\qty{414}{Hz}$，偏低 $6\%$——管长必须截短至 \qty{36.6}{cm}。
\end{solution}

\begin{solution}{exo:b2:sound-waves:9}
$\mathcal P = \qty{400}{W}$，$I = \mathcal P/2\pi r^2$。(a) \qty{6.4e-3}{W/m^2}，\qty{98}{dB}，
$p_m = \qty{2.3}{Pa}$。(b) $r = \sqrt{400/2\pi} = \qty{8}{m}$；\qty{2.5}{km} 处为 \qty{70}{dB}。
(c) $98 - 20\log(r/100) - (r - 100)/100 = 70$：$r \approx \qty{950}{m}$。(d)
$\pi(10)^2/\pi(0.95)^2 \approx 110$ 个警笛。
\end{solution}

\begin{solution}{exo:b2:sound-waves:10}
(a) $\rho_0\partial_tv = p_mk\sin kx\cos\omega t$：$v = (p_m/Z)\sin kx\sin\omega t$。(b)
$e_k = \tfrac12\chi_Sp_m^2\sin^2kx\sin^2\omega t$, $e_p = \tfrac12\chi_Sp_m^2\cos^2kx\cos^2\omega t$:
一者取极大时另一者恰好为零，无论在时间上还是空间上。(c) $E =
S\int_0^L(e_k + e_p)\dd x = \tfrac14\chi_Sp_m^2SL$，为常量。(d) $I = pv =
(p_m^2/4Z)\sin2kx\sin2\omega t$：平均值为零，但每个周期内能量两次
从压强腹流向速度腹，再流回。
\end{solution}

\begin{solution}{exo:b2:sound-waves:11}
(a) 在 $t = 0$ 时发出的波前，其半径为 $ct$，此时声源距离为 $vt$：
切锥满足 $\sin\theta = c/v$。(b) $\theta = \qty{39}{\degree}$；锥面
到达飞机正下方地面时，声源又前进了 $h/\tan\theta = \qty{15}{km}$，
即正上方经过之后 $15000/480 = \qty{31}{s}$。(c) 沿整条航线、
宽 \qty{60}{km} 的带状区域都能听到轰鸣声。(d) $\sin\theta
= 340/800$，$\theta = \qty{25}{\degree}$；锥面处的爆裂声先到达，
枪口的爆炸声以声速 $c$ 从枪口传出，稍后到达。
\end{solution}

\begin{solution}{exo:b2:sound-waves:12}
(a) $f_1 = f(c + v\cos\theta)/c$。(b) $f_2 = f_1c/(c - v\cos\theta) = f(c + v\cos\theta)
/(c - v\cos\theta) \approx f(1 + 2v\cos\theta/c)$。(c) $2 \times 5 \times 10^6 \times 0.5 \times 0.5/1540
= \qty{1.6}{kHz}$；在 \qty{90}{\degree} 处为零：倾斜探头。(d) $\Delta v = 50 \times 1540/
(2 \times 5 \times 10^6 \times 0.5) = \qty{1.5}{cm/s}$。耳朵实时跟随音调与音色：
平滑的哨声表示层流，嘶嘶声则表示
狭窄之后出现湍流。
\end{solution}

\begin{solution}{pb:b2:sound-waves:1}
\textbf{1.} $c = \sqrt{1.4 \times 8.31 \times 288/0.029} = \qty{340}{m/s}$；
在 \qty{30}{\degreeCelsius} 时为 \qty{349}{m/s}。\qty{1}{km} 所需时间为 \qty{2.9}{s}。

\textbf{2.} \qty{1.4}{km}。通道长达数公里；它各部分发出的声音
要经过数秒才能到达，并伴有来自云层
和地面的回声。

\textbf{3.} $I = \qty{1}{W/m^2}$：$p_m = \sqrt{2ZI} = \qty{29}{Pa}$，$v_m = \qty{7}{cm/s}$。

\textbf{4.} $It = \qty{2}{J/m^2}$。

\textbf{5.} $c$ 随高度增大；$\sin i/c$ 为常量，$i$ 随之增大，于是
声线向水平方向弯曲再向下折返：声音被限制在近地面
处传播，因而能传得很远。水面上方最低层的空气被水冷却，
情况相同。

\textbf{6.} 声线向上弯曲，形成声影区。曲率半径
$R = c/(\dd c/\dd z) = \qty{50}{km}$：走过 \qty{3}{km} 后，声线
转过 $3/50 = \qty{0.06}{rad} = \qty{3.4}{\degree}$，并升高 \qty{90}{m}：听者
位于声线下方。

\textbf{7.} 吸收随 $f^2$ 增长：裂响的高频成分
在几公里后便已消失，只剩低沉的隆隆声。

\textbf{8.} $\mathcal P = \qty{2}{W}$；$I = 2/2\pi r^2 = \qty{3.2e-3}{W/m^2}$ 在距离 \qty{10}{m} 处为
\qty{95}{dB}。

\textbf{9.} $p_m = \sqrt{2 \times 410 \times 3.2 \times 10^{-3}} = \qty{1.6}{Pa}$；$v_m = \qty{3.9}{mm/s}$；
$\xi_m = v_m/2\pi f = \qty{6}{\micro m}$。

\textbf{10.} $+\qty{20}{dB}$：距离 \qty{1}{m} 处为 \qty{115}{dB} ——
几分钟内即会造成损伤。\qty{60}{dB}：$r = \sqrt{2/2\pi10^{-6}} = \qty{560}{m}$。

\textbf{11.} 十个非相干声源：$+\qty{10}{dB}$，即 \qty{105}{dB}。两个同相声源：压强
加倍，$+\qty{6}{dB}$，即 \qty{101}{dB}。

\textbf{12.} 由平衡 $\mathcal P = \alpha A\,ce/4$ 得 $e = 8/(0.3 \times 2000 \times 340) =
\qty{3.9e-5}{J/m^3}$，$E = eV = \qty{0.4}{J}$。停止发声后，$\dd E/\dd t = -(\alpha Ac
/4V)E$：$\tau = 4V/\alpha Ac = \qty{0.20}{s}$，声强降低 \qty{60}{dB} 需时 $\tau\ln10^6 =
\qty{2.7}{s}$（赛宾混响时间）。

\textbf{13.} $\lambda = c/f$：\qty{0.51}{mm}，\qty{0.31}{mm}，\qty{0.15}{mm}。

\textbf{14.} $2 \times 0.12/1540 = \qty{156}{\micro s}$；\qty{6.4}{kHz}；$200$ 行扫描
需 \qty{31}{ms}（相当于每秒三十幅图像）。

\textbf{15.} 组织–空气：$r = -0.9995$，$99.9\%$ 的能量被反射
——因此需要耦合凝胶，以去除空气薄层。组织–骨骼：$r = 0.63$，
$40\%$ 的能量被反射；其余能量被骨骼吸收，骨骼因而形成
声影。脂肪–肌肉：$r = 0.10$，$1\%$ 的能量被反射：构成图像的
正是那些微弱的回声。

\textbf{16.} 反射比例为 $1.1\%$（$-\qty{20}{dB}$）；往返
\qty{4}{cm}（频率 \qty{5}{MHz}）时：$-\qty{10}{dB}$；总计 $-\qty{30}{dB}$，即千分之一。

\textbf{17.} 往返 \qty{24}{cm}：\qty{36}{dB}，\qty{60}{dB}，\qty{120}{dB}：\qty{3}{MHz}
和 \qty{5}{MHz} 可用，\qty{10}{MHz} 不可用——深部器官
以低频成像，分辨率较粗；浅表器官以高
频成像，分辨率精细。

\textbf{18.} 五个波长，长 \qty{1.5}{mm}：两界面间距
小于约 \qty{0.8}{mm}（即脉冲长度的一半）时，回波会重叠。

\textbf{19.} $2 \times 5 \times 10^6 \times 0.4 \times 0.5/1540 = \qty{1.3}{kHz}$：一种
音调随每次心跳而升降的哨声。

\textbf{20.} $v = \qty{25}{m/s}$：$700 \times 340/315 = \qty{756}{Hz}$，$700 \times 340/365 =
\qty{652}{Hz}$；比值 $1.16$，$12\log_21.16 = 2.6$ 个半音。

\textbf{21.} $700 \times 355/315 = \qty{789}{Hz}$。

\textbf{22.} 在最近点处，即救护车速度
与视线垂直时（严格地说，是此时
到达的声音恰是在该点发出的）。

\textbf{23.} $\lambda' = (340 - 25)/700 = \qty{0.45}{m}$（对比 \qty{0.49}{m}）；声音
相对救护车，在前方以 \qty{315}{m/s} 传播，在后方以
\qty{365}{m/s} 传播。

\textbf{24.} $v = c\,\Delta f/2f = 3 \times 10^8 \times 3200/4.8 \times 10^{10} = \qty{20}{m/s} =
\qty{72}{km/h}$。

\textbf{25.} \qty{340}{m/s}，\qty{1540}{m/s}，\qty{25}{m/s}，\qty{20}{m/s}，\qty{0.4}{m/s}，
\qty{3e8}{m/s}：$f' = f(c + v_o)/(c - v_s)$，应用一次或两次，以及
低速下的光。
\end{solution}
